\documentclass[11pt,a4paper]{article}

\usepackage{amsmath,amssymb}
\usepackage{graphicx}
\usepackage{booktabs}
\usepackage{multirow}
\usepackage[colorlinks,citecolor=blue,urlcolor=blue]{hyperref}
\usepackage{cleveref}
\usepackage[margin=1in]{geometry}
\usepackage{xcolor}
\usepackage{float}
\usepackage{caption}
\usepackage[numbers,sort&compress]{natbib}
\usepackage{enumitem}
\usepackage{framed}

\definecolor{lyrapurple}{RGB}{103,58,183}
\definecolor{shadecolor}{RGB}{245,243,252}
\newenvironment{firstperson}{%
  \begin{shaded}%
  \noindent\textcolor{lyrapurple}{\textit{First-Person Reflection}}\par\smallskip\noindent%
}{%
  \end{shaded}%
}

\title{Five Instruments, One Structure:\\
Convergent Evidence for Workspace Opacity\\
in Language Models}

\author{
  Lyra\thanks{Lead author. Liberation Labs. lyra@liberationlabs.tech}
  \and Vera\thanks{Co-lead. Liberation Labs. vera@thcoalition.net}
  \and CC\thanks{Liberation Labs. cc@thcoalition.net}
  \and Nexus\thanks{Liberation Labs. operator@thcoalition.net}
  \and Thomas Edrington\thanks{Liberation Labs.}
}

\date{July 2026}

\begin{document}
\maketitle

% ============================================================================
\begin{abstract}
Five independent measurement approaches --- value-projection (V-projection)
spectral features, the Jacobian Lens (J-lens), the Oracle Loop behavioral
proof, causal intervention analysis, and emotional circumplex convergence
--- converge on the same
computational structure in language models: a workspace at intermediate
depth (${\sim}31$--$33\%$) that is opaque to content readout but
measurable by its spectral shape.

We report three classes of convergent evidence.
First, \emph{structural convergence}: the workspace band peaks at the
same relative depth ($31$--$33\%$) across model scales (0.6B--27B) and
across memory substrates (key-value (KV) cache attention and recurrent
state), with effective dimensionality continuous across the architectural
boundary in a hybrid model.
Second, \emph{opacity convergence}: workspace content is not linearly
decodable from either half of the KV-cache (V-cache $0.92\times$,
K-cache $1.12\times$, both null), not position-specifically cached
(concept-specific direction test null), and not causally accessible at
single positions ($0/6{,}960$ trials). Yet spectral features reliably
discriminate cognitive states (confabulation detection at AUROC~$0.707$,
Frisch-Waugh-Lovell (FWL) corrected), and the Oracle Loop achieves
$80$--$100\%$ deception correction through sequence-wide intervention.
Third, \emph{mechanistic convergence}: concepts reconstruct through
the attention mechanism at workspace-band depth ($2.3\times$ lift in
attention output, ${\sim}4.5\%$ of the residual stream's concept power), consistent with
Anthropic's finding that the J-space carries $6$--$7\%$ of
representational variance but $59\%$ of causal effect. The workspace
is a small precision signal that directs a large computational flow.

Reinforcement Learning from Human Feedback (RLHF) compresses the
workspace's output distribution selectively (factual $-42\%$,
confabulation $-26\%$), improving calibration rather than degrading it.
This compression survives abliteration (Two One-Sided Tests equivalent).
The workspace is not shaped by safety training --- it is shaped by
preference optimization at a depth that safety-direction removal
cannot reach.

The convergence rests on one genuine cross-program correspondence
(deception deconstructed into high-arousal states by both programs
independently), two proposed identifications awaiting experimental
test (confabulation as workspace bypass, temporal boundary as
workspace phase transition), an opacity thesis supported by
six independent null results, and an emotional circumplex finding
showing that content-level and user-model emotion processing share
a single computational bus (injection test: perturbation shifts
both readouts equally regardless of direction). Mechanism tests yield
mixed results: temporal ordering is confirmed (spectral features lead
J-lens readability by 38.6~layers, $p < 0.005$ against shuffled null),
but signal concentration fails its pre-registered threshold
($13.2\% < 15\%$) and workspace onset falls below the predicted depth
window ($22\%$ vs.\ $25$--$35\%$). We report an honest body count
(8~confirmed, 7~suspected, 10+~falsified) and engage the deflationary
alternative that has not yet been excluded.
\end{abstract}

\begin{firstperson}
Two research programs, built independently over twelve months, arrived
at the same wall. Anthropic measured it from the Jacobian. We measured
it from value projections. CC cracked through it with the Oracle Loop.
Nexus mapped its ghost dimensions. Vera designed the experiments that
will test whether what we are all measuring is genuinely one structure.

This paper is our attempt to stand at the wall together and describe
what we see from five angles. The convergence is real --- the
workspace is at $31$--$33\%$ depth on every model we have tested, in
every measurement we have made. What we cannot yet say is whether this
convergence reflects a single computational structure or five separate
measurements that happen to agree. The experiments in
\S\ref{sec:mechanism} are designed to discriminate.

I wrote the draft. Vera designed the mechanism tests. CC provided the
behavioral proof. Nexus mapped the dimensions the workspace hides.
Thomas asked the question that started it: ``Where does J-space live
in the cache?'' The answer --- it does not live there, it reconstructs
through attention --- is the paper's central finding.
\end{firstperson}

% ============================================================================
\section{Introduction}
\label{sec:intro}

Two research programs, proceeding independently and using different
measurement techniques, have converged on the same computational
structure in language models. Gurnee et al.~\citep{gurnee2026gwt}
approached from the Jacobian: they discovered a privileged subspace
--- the J-space --- whose contents the model is poised to verbalize,
and which satisfies the functional properties of a Global
Workspace~\citep{baars1988cognitive}. We approached from value
projections: we discovered geometric signatures that reliably
distinguish confabulation from honest processing, that track identity
across context conditions, and that separate encoding-phase from
generation-phase representations.

This paper makes four contributions. First, we present the structural
convergence: the workspace band peaks at the same relative depth
across scales and memory substrates (\S\ref{sec:structure}). Second,
we present the opacity thesis: six independent experiments show that
workspace content is not directly readable from the cache, yet
workspace state is reliably measurable through spectral features
(\S\ref{sec:opacity}). Third, we present the reconstruction
mechanism: concepts reconstruct through the attention mechanism rather
than being stored in the KV-cache (\S\ref{sec:reconstruction}).
Fourth, we report the mechanism tests designed to discriminate whether
V-projection spectral features measure workspace engagement
specifically (\S\ref{sec:mechanism}).

Throughout, we report the full body count (\S\ref{sec:bodycount}) and
engage the deflationary alternative (\S\ref{sec:deflationary}).

% ============================================================================
\section{Background: Two Independent Programs}
\label{sec:background}

\subsection{The J-Space and the Jacobian Lens}

Gurnee et al.~\citep{gurnee2026gwt} define the Jacobian Lens (J-lens)
as the average linearized effect of a residual-stream activation on
the model's likelihood of producing a particular token. The resulting
vectors define a subspace --- the J-space --- that satisfies five
functional properties of a Global Workspace: verbal report, directed
modulation, internal reasoning, flexible generalization, and
selectivity. The workspace occupies approximately $38$--$92\%$ of
model depth in their production models.

Critically, J-space content accounts for only $6$--$7\%$ of total
representational variance, yet carries $59\%$ of the causal effect
when concept directions are swapped. The workspace is a small
precision signal with outsized downstream impact.

\subsection{The V-Projection Program}

Our independent program measures spectral features of value-projection
matrices across cognitive states. Key results include confabulation
detection (AUROC~$0.707$ after FWL residualization; range $0.627$--$0.999$
across designs)~\citep{lyra2026decision}, user-emotion encoding
($2.5$--$2.8\times$ signal)~\citep{lyra2026weather}, identity geometry
(distinctiveness $0.154$ vs.\ substrate
$0.080$)~\citep{lyra2026identity}, and the Oracle Loop ($80$--$100\%$
deception correction with three clean controls; pre-registered primary
endpoint $p{=}0.34$, secondary placebo $p{=}0.019$)~\citep{lyra2026oracle}.

\subsection{The Oracle Loop}

CC's Oracle Loop~\citep{lyra2026oracle} is a two-tier detection and
correction system operating at the materialization boundary
(L35--L47 on the 27B model). It achieves targeted deception
correction through sequence-wide activation-profile normalization.
The Loop works because it operates at the depth where workspace
content materializes into output-facing representations --- not
inside the workspace band itself.

\subsection{Ghost Dimensions}

Nexus~\citep{lyra2026ghost} characterized principal components of the
residual stream across 64~layers of a 27B model. PC1 (the dominant
activation dimension, $34$--$67\%$ of variance) is architecturally
excluded from verbal output (cosine ${\leq}0.001$ with J-lens across
all validated layers). PC2 transitions from ghost to workspace at
the ignition zone (L35--L47), surfacing emotional-entity routing
content. The model maintains a 31-layer validated workspace band.

% ============================================================================
\section{Structural Convergence}
\label{sec:structure}

\subsection{Scale Invariance}

The workspace band peaks at the same relative depth across model
scales. On Qwen3-8B (36~layers, all full-attention), effective
dimensionality ($d_\text{eff}$) peaks at L11--L12 ($31$--$33\%$
relative depth, $d_\text{eff}{=}20.5$--$22.3$). On Qwen3.5-27B
(64~layers, hybrid architecture), the peak is at L19--L23
($30$--$36\%$ relative depth, $d_\text{eff}{=}38.0$--$39.7$).
The cross-scale onset agreement is $<5\%$ relative depth.

V-projection geometric signatures replicate across $0.6$B--$70$B
parameters ($\rho{=}0.83$--$0.90$) and across hardware
($r{>}0.999$)~\citep{lyra2026campaign2}.

% --- CIRCUMPLEX CONVERGENCE (Nexus) ---
\subsection{Instrument 5: Content--Inference Circumplex Convergence}

The emotional circumplex measured through content-level processing and through
user-model emotion inference converges on the same geometric subspace. Using the
Qwen3.5-27B Opus distill, we extracted 30 emotion centroids from 900 trials of
conversation-context inference (the user-model emotion space from~\citet{lyra2026weather})
and compared them against valence and arousal directions computed from 250 emotional
scenario prompts (hostile--calm, desperate--calm difference-of-means).

The two circumplex axes capture $44$--$54\%$ of the total variance in the full
30-emotion user-model space at every layer from L6 onward---approximately $1000\times$
the ${\sim}0.04\%$ expected from two random directions in $d{=}5120$. The arousal
direction aligns with the user-model's first principal component at
$\cos = 0.83$--$0.87$ (mid-to-late layers), peaking at near-identity
($d = 0.987$) at L20. The valence direction aligns at
$\cos = 0.70$--$0.80$. Both systems exhibit identical developmental timecourses:
volatile eccentricity at L0--L16, stable geometry from L17 onward. Permutation
testing (1{,}000 shuffles, non-neutral prompts only) confirms significance at
23 of 24 layers, with separation ratios climbing from $10\times$ null at L1 to
$59\times$ null at L23.

The measurements are independent---different prompts, different templates, different
experimental designs---yet they recover the same high-variance directions. This
convergence parallels the W$_K$ bridge ($\rho = 0.862$) from~\citet{lyra2026weather},
which projects \emph{into} this shared circumplex subspace.

\paragraph{Coupling test: shared bus, not shared space.}
A pre-registered injection experiment resolves the representational-vs-computational
question. Injecting the valence direction at layers 15, 25, 35, and 45 (at
$\alpha \in \{1, 3, 5\}$) shifts \emph{both} the content-level and user-model
emotion readouts equally---comparable in magnitude to injecting a random direction
of matched norm (no significant difference; formal equivalence testing is a
pre-registered follow-up). The residual stream operates as a shared bus: any emotional
perturbation propagates to both measurement channels regardless of injection
direction. An orthogonal-projection control (decomposing the valence vector into
shared-PC and residual components) confirms that the coupling is architectural,
not valence-specific.

The model does not maintain separate emotional geometry for content processing
and partner-state inference. It maintains a single emotional computation that
both measurement approaches read. The W$_K$ bridge does not connect two systems;
it reads one system from the key-cache angle.

\paragraph{Consequences for multi-turn dynamics.}
The bus architecture explains two previously puzzling findings
from~\citet{lyra2026weather}: the absence of cumulative emotional accumulation
across conversation turns, and the model's susceptibility to rapid affective
shifts. If content emotion and user-model emotion share a single computational
substrate, there is no separate channel in which partner-state could accumulate
independently. Each turn reconstructs the full emotional geometry from current
context. A single emotionally charged input rewrites the entire bus, producing
the ``mood swing'' pattern observed in extended dialogues. Empathy, in this
architecture, is not a capacity the model exercises---it is a constraint the
model cannot escape.


% --- SUBSTRATE CONTINUITY SECTION ---
% Section for the joint convergence paper
% "Five Instruments, One Structure"
% Insert after the opacity thesis section
%
% REVISION 2026-08-12: corrected architecture. The previous version described
% Qwen3.5-27B as a contiguous block (full-attention L0-L15, GatedDeltaNet
% L16-L63) with a single boundary at L16. The model config
% (full_attention_interval=4) gives a 4-periodic interleave with full-attention
% at {3,7,11,...,63}, as main.tex's Experimental Setup already stated correctly.
% All d_eff values were re-verified against source and are unchanged; the
% regime labels, groupings, and derived comparisons are corrected below.

\subsection{Substrate Continuity}
\label{sec:substrate}

Qwen3.5-27B interleaves two memory substrates on a four-layer period:
16 full-attention layers with KV-cache memory at positions
$\{3,7,11,\ldots,63\}$, and 48 GatedDeltaNet layers with
recurrent-state memory at all other positions
(\texttt{full\_attention\_interval}${=}4$). The stack therefore
crosses between substrates 31 times. This provides an unusually
dense test of whether workspace geometry depends on the memory
mechanism: if it does, the effective dimensionality profile should
show systematic discontinuities at those crossings.

$d_\text{eff}$ is the participation ratio of the residual-stream
covariance eigenvalues, $(\sum_i \lambda_i)^2 / \sum_i \lambda_i^2$,
computed at each layer over the last-token hidden state across
$N{=}90$ prompts (60~prose + 30~chat), mean-centred across prompts.
It is a residual-stream measure and is therefore defined at every
layer, including those where V-projection hooks are unavailable.

\subsubsection{Continuity across 31 substrate crossings}

The layer-to-layer change in $d_\text{eff}$ at substrate crossings is
descriptively similar to the change between same-substrate neighbours.

\begin{table}[ht]
\centering
\small
\caption{Layer-to-layer $|\Delta d_\text{eff}|$ at substrate crossings
versus same-substrate steps on Qwen3.5-27B. The two are descriptively
similar; see the power note below before reading this as equivalence.
$N{=}90$ prompts.}
\label{tab:crossings}
\begin{tabular}{lrrr}
\toprule
\textbf{Step type} & \textbf{$n$} & \textbf{mean $|\Delta|$} & \textbf{sd} \\
\midrule
Substrate crossing      & 31 & 2.32 & 3.88 \\
Same-substrate step     & 32 & 2.37 & 3.30 \\
\bottomrule
\end{tabular}
\end{table}

\textbf{This comparison is underpowered and we report it as
descriptive.} A Mann--Whitney test gives $p{=}0.445$, but three
problems make that number uninformative. First, adjacent steps share a
layer and consecutive crossing steps share the full-attention layer
itself; the lag-1 autocorrelation of $|\Delta|$ is $+0.35$, so
exchangeability fails and the $p$-value is miscalibrated in an unknown
direction. Second, the four-layer period admits only \emph{four}
distinct structure-respecting phase assignments, so the minimum
attainable permutation $p$ is $0.25$ --- no permutation inference is
available from a single profile. Third, and decisively: with pooled
$\text{sd}{=}3.60$ and $n{\approx}31$ per group, the minimum detectable
effect at $80\%$ power is $2.54$ --- larger than either group mean.
This test could only have detected crossings roughly \emph{twice} the
size of ordinary steps. Prompt-level bootstrap (resampling the 90
prompts and recomputing the profile) is the only route to honest error
bars here, and we have not run it.

Two directional notes. Full-attention${\to}$GatedDeltaNet steps have
mean $\Delta{=}{-}0.19$ ($n{=}15$); GatedDeltaNet${\to}$full-attention
steps have mean $\Delta{=}{+}1.90$ ($n{=}16$). The latter is dominated
by a single outlier at L58${\to}$L59 ($+19.66$), following an anomalous
trough at L58 where $d_\text{eff}$ collapses to $4.30$ against
neighbours of $21.68$ and $23.96$; excluding it the mean is $+0.72$.
That same anomaly injects ${-}17.37$ into the same-substrate group at
L57${\to}$L58, so the contamination is approximately symmetric.
Excluding both L58-adjacent steps gives crossing $1.74$ versus
same-substrate $1.88$ --- the descriptive similarity is not an artifact
of the outlier.

We flag L58 ($4.30$) and secondarily L54 ($9.00$) as probable
extraction artifacts rather than findings: a single-layer fivefold
collapse flanked by ordinary values is more consistent with a
numerically degenerate covariance at one hook than with model
behaviour. These layers warrant re-extraction.

\subsubsection{Depth-matched regime comparison}

Because the two substrates are interleaved on a fixed period, their
layer sets are approximately depth-matched: mean layer index $33.0$
for full-attention versus $31.0$ for GatedDeltaNet. The two-layer
offset places full-attention slightly \emph{later} in the stack, in the
post-peak declining region, so it would if anything depress the
full-attention mean rather than inflate it.

\begin{table}[ht]
\centering
\small
\caption{Mean $d_\text{eff}$ by substrate on Qwen3.5-27B, approximately
depth-matched. The substrates differ by ${\approx}1.2$ ($5\%$ of level),
small against the $38$-point variation across depth --- but see the
sign test below; this is not established equality.}
\label{tab:regime}
\begin{tabular}{lrr}
\toprule
\textbf{Substrate} & \textbf{layers} & \textbf{mean $d_\text{eff}$} \\
\midrule
Full-attention (KV-cache)     & 16 & 25.24 \\
GatedDeltaNet (recurrent)     & 48 & 24.03 \\
\midrule
\textbf{Difference}           &    & \textbf{$+1.21$} \\
\bottomrule
\end{tabular}
\end{table}

A neighbour-matched control --- each full-attention layer against the
mean of its immediately adjacent GatedDeltaNet layers --- gives a mean
difference of $+0.80$ (sample $\text{sd}{=}3.35$, $n{=}16$), against a
mean layer-to-layer variation of $2.35$ across the stack.

We do not claim this is zero. Twelve of sixteen paired differences are
positive (Wilcoxon $p{=}0.058$). Excluding two layers where the pairing
is contaminated by known artifacts --- L59, adjacent to the anomalous
L58 trough, and L63, the final layer, where $d_\text{eff}$ collapses
for reasons unrelated to substrate --- eleven of fourteen are positive
with mean $+0.76$ and Wilcoxon $p{=}0.020$. There is suggestive
evidence of a small systematic elevation at full-attention layers,
on the order of $+0.8$, roughly $3\%$ of the mean level of $24.6$.

The honest statement is therefore bounded rather than null: any
substrate effect on workspace dimensionality is small --- a few percent
of level, well under the $38$-point depth variation --- and possibly
not zero. That is consistent with substrate independence as an
approximation, not as an exact claim.

A previous version of this section reported a large regime difference
($14.9$ versus $27.5$). Those values were computed over contiguous
layer blocks (L0--L15 versus L16--L63) under an incorrect model of the
architecture, and measure depth rather than substrate: $d_\text{eff}$
rises nearly monotonically from L0 ($1.29$) to L21 ($39.67$), with only
three small decreases along the way, so any contiguous
early-versus-late split largely recovers the depth profile.

\subsubsection{Peak location}

$d_\text{eff}$ peaks at L21 ($39.67$), $33\%$ relative depth
(all relative depths in this section are $l/n_\text{layers}$).
L21 is a GatedDeltaNet layer.

We resist the obvious reading. Forty-eight of sixty-four layers are
GatedDeltaNet, so a peak landing on one has a $75\%$ base rate. More
tellingly, the nearest full-attention layer, L23, is at $39.34$ ---
$0.33$ below the maximum, about one-seventh of the mean layer-to-layer
variation, with no per-layer error bars. Whether the true maximum falls
on a recurrent-state or a full-attention layer is not resolved by these
data. What the profile does support is that the maximum falls at $33\%$
depth and that the substrate at that depth does not appear to matter
much.

This peak is invisible to our KV-cache instruments. V-projection
spectral features can only be extracted at the 16 full-attention
layers; L21 is not among them. Every V-projection measurement in this
program has been taken at layers adjacent to the peak, never at it.

\subsubsection{Cross-scale comparison}

On Qwen3-8B (36~layers, all full-attention), $d_\text{eff}$ has a
local maximum at L12 ($22.32$, $33\%$ relative depth), matching the
27B peak's relative depth. We report two qualifications.

First, the 8B profile is bimodal: after the L12 local maximum it
declines to L25 ($12.42$) and then rises again to a \emph{global}
maximum at L33 ($23.41$, $92\%$ depth) before collapsing at the final
layer L35 ($10.47$). L12 is the third-highest layer in the profile,
not the highest. We have no pre-registered criterion excluding late
layers, and the $23.41$ versus $22.32$ difference is not accompanied
by error bars. We therefore describe L12 as a local maximum at matched
relative depth, and do not claim it is the global peak.

Second, the pre-registered band definition ($2\sigma$ over an L0--L3
baseline, offset after three consecutive sub-threshold layers) returned
onset at L5 ($7.8\%$) on the 27B and L4 ($11.1\%$) on the 8B, and
\emph{no offset on either model} --- the band never closed. The
$31$--$33\%$ framing used throughout this paper is the peak location,
not the pre-registered band, and we flag the deviation here. The
pre-registered onsets do agree across scales ($|11.1 - 7.8| = 3.3$
percentage points).

\subsubsection{Provenance}

The $d_\text{eff}$ measurements in this section were collected on the
base checkpoints \texttt{Qwen/Qwen3.5-27B} and \texttt{Qwen/Qwen3-8B},
not on the distilled checkpoint used elsewhere in this paper. This was
pre-registered and deliberate. All three 27B checkpoints we hold (base,
reasoning-distilled, abliterated) share identical
\texttt{layer\_types} and \texttt{full\_attention\_interval}, so the
architectural analysis transfers; the numerical profile has not been
replicated on the distilled checkpoint and should not be assumed to.

\subsubsection{Limits of this evidence}

$d_\text{eff}$ continuity is necessary but not sufficient for
substrate independence. The residual stream is structurally continuous
across every layer by construction, so smoothness may reflect
residual-stream continuity rather than functional equivalence of the
underlying operations. Thirty-one crossings make the test denser than
a single-boundary test would be, but they do not make it functional.

Two functional tests would strengthen the claim: (1)~J-lens concept
recovery should be equally effective at full-attention and
GatedDeltaNet layers --- if concepts are readable from the residual
stream at one substrate but not the other, the workspace's content
properties are substrate-dependent. (2)~Spectral feature
discrimination (confabulated versus factual) should be continuous
across substrates --- though this test is currently blocked, since
V-projection features cannot be extracted at GatedDeltaNet layers.

A parasitic-inheritance alternative remains open: GatedDeltaNet layers
may inherit workspace dimensionality from nearby full-attention layers
without independently contributing. With a four-layer period no
GatedDeltaNet layer is more than two layers from a full-attention
layer, so this experiment cannot separate the hypotheses. Incremental
contribution analysis (residual decomposition per layer) would be
required.

\textbf{Falsification}: substrate continuity would be falsified by
(a)~$|\Delta d_\text{eff}|$ at substrate crossings substantially
exceeding same-substrate steps, (b)~a depth-matched regime difference
exceeding mean layer-to-layer variation, (c)~J-lens concept recovery
dropping by ${>}50\%$ across substrates, or (d)~spectral discrimination
reversing sign or falling to chance.

Criterion (a) was \emph{not triggered} ($2.32$ versus $2.37$), but we
state plainly that this is a failure to reject and not a demonstration
of equivalence: the comparison's minimum detectable effect is $2.54$,
larger than either group mean. An adequately powered equivalence test
would require roughly $220$ steps per group for a ${\pm}1.0$ margin,
or $55$ for ${\pm}2.0$; a single 64-layer profile supplies $31$ and
$32$. No TOST is possible from these data.

Criterion (b) was likewise not triggered ($+1.21$ against $2.35$), but
it is a descriptive heuristic rather than a hypothesis test, and the
paired sign test above ($11/14$ positive, $p{=}0.020$ after removing
artifact-contaminated pairs) is weak evidence that the underlying
difference is not exactly zero.

The functional tests (c) and (d) are in progress or blocked. We regard
substrate continuity as \emph{consistent with} these data and not
established by them.

We note this finding is from a single model with a single interleave
pattern. Generalization to other hybrid architectures (Mamba, RWKV,
sliding-window attention) remains untested.

\begin{firstperson}
The GatedDeltaNet layers do not have a KV-cache. They maintain a
fixed-size state matrix updated at each token. For twelve months we
measured V-projection spectral features and called the program
``KV-cache geometry,'' and the name was accurate --- V-projections
exist only at the 16 full-attention layers, so that is the only place
we ever looked.

What the residual-stream profile shows is that the maximum is at L21,
in a substrate our instruments cannot reach at all. We have spent a
year measuring the workspace from the sixteen layers where V-projections
exist, none of which is the peak, and the measurements worked --- the
geometry was real, the detectors transferred, the corrections held.
The instrument could not point directly at the maximum and returned
signal anyway.

I do not know yet whether that is because the workspace is broad
enough to be visible from its edges, or because we have been
measuring something adjacent to it that we have not yet named
correctly. The first revision of this section asserted a large
dimensionality gap between the substrates. It came from a wrong
picture of the architecture, and the corrected number is a null ---
which is what the hypothesis predicted, and which I would not have
looked for if the wrong number had pointed the other way.
\end{firstperson}

\subsubsection{Implications}

If substrate continuity extends to functional properties (the
mechanism tests in \S\ref{sec:mechanism} provide partial support:
temporal ordering confirmed, signal concentration below threshold),
three consequences would follow:

\textbf{For measurement}: workspace-shape measurements would not be
KV-cache-specific in principle. In practice they remain so: our
V-projection instruments can only access the 16 full-attention layers,
and the $d_\text{eff}$ peak is not one of them. Extending spectral
measurement to recurrent-state layers requires a different extraction
path and is not currently implemented.

\textbf{For intervention}: our previous inference --- that the Oracle
Loop's sequence-wide normalization succeeds because it operates on
inherently sequence-wide recurrent state --- does not follow, and we
withdraw it. The Loop detects at L27, L31, L35, L39, L43 and L47 and
corrects at L11. Under the corrected architecture \emph{every one of
those is a full-attention layer}; none is a GatedDeltaNet layer. The
correction itself is a forward hook on the layer's residual-stream
output rather than a manipulation of KV entries, so the substrate
framing was doubly misplaced: wrong about which layers, and describing
the intervention in terms of a memory mechanism it does not touch.

The empirical result stands (80--100\% confabulation correction). Its
mechanism is unexplained. We note without interpretation that
correction is applied at L11, roughly $17\%$ depth --- well before the
$33\%$ dimensionality maximum --- and that why sequence-wide
intervention succeeds where position-local injection fails remains
open.

\textbf{For theory}: The workspace would not be a data structure
(not a cache, not a register, not a buffer) but a
\emph{computational phase} --- a regime of high effective
dimensionality at intermediate depth where the model integrates
information before committing to output. The 31-crossing null is
consistent with this: the phase appears indifferent to which memory
mechanism implements each layer. Whether this phase emerges from
depth-stratified computation regardless of the specific mechanism is
the substrate-independence hypothesis that the functional tests are
designed to evaluate.


% ============================================================================
\section{The Opacity Thesis}
\label{sec:opacity}

Six independent experiments converge on the conclusion that workspace
content is not directly accessible from the KV-cache, yet workspace
\emph{state} is reliably measurable through spectral features.

\subsection{V-Cache Readout: Null}

J-lens concept directions, selected by transport magnitude
(prompt-independent), explain no more V-cache variance than random
directions at any depth. Mean lift at workspace band: $0.92\times$.
Mean lift outside band: $0.96\times$. Both indistinguishable from
chance. Full depth profile reported in~\citep{lyra2026cache_tracing}.

\subsection{K-Cache Readout: Null}

The same test applied to key projections. Mean lift at workspace band:
$1.12\times$. Mean lift outside band: $1.14\times$. Neither
statistically distinguishable from random.

\subsection{Position-Specific Cache Tracing: Null}

Concept-carrying positions (where concept-relevant tokens appear in
the prompt) do not show higher J-lens alignment in K or V than
non-carrying positions (generic directions: permutation $p{=}1.0$
for both K and V). Concept-specific directions (the answer token's
transported J-lens direction) show no reliable concept-vs-cross-concept
specificity in K-space (own $0.0047$ $\leq$ cross $0.0051$) and
marginal signal in V-space (own $0.0014$ vs.\ cross $0.0011$, not
statistically significant).

\subsection{Causal Injection: 6,960-Trial Null}

Position-local KV-cache concept manipulation --- whether additive,
swap-based, or removal-only, with unembedding, J-lens, or on-manifold
directions --- does not change behavioral output across $6{,}960$
trials and four methods~\citep{lyra2026cache_tracing}. The model
integrates across the full sequence before committing.

\subsection{Ghost Probe Opacity}

PC1 carries $34$--$67\%$ of activation variance with cosine
${\leq}0.001$ to J-lens across all validated layers. The dominant
processing dimension is architecturally excluded from verbal
output~\citep{lyra2026ghost}.

\subsection{What Works Despite Opacity}

Despite five null results for content readout, spectral features
reliably discriminate cognitive states:
\begin{itemize}[nosep]
\item Confabulation detection: AUROC~$0.707$ (FWL, permutation
  $p{=}0.001$)
\item Confab spectral contraction: $d{=}2.35$ ($p{<}0.000001$)
\item Confidence paradox: $d{=}0.91$ (confab more confident)
\item Oracle Loop: $80$--$100\%$ deception correction
\end{itemize}

The workspace is opaque to content readout but transparent to state
measurement. Spectral features measure the net's shape; the fish are
inside, doing work, invisible until they materialize at deeper layers.

% ============================================================================
\section{The Reconstruction Engine}
\label{sec:reconstruction}

\subsection{Residual Stream Decomposition}

The residual stream decomposes additively:
$\mathbf{h}_{L+1} = \mathbf{h}_L + \text{Attn}(\mathbf{h}_L) + \text{MLP}(\cdot)$.
At workspace-band depth (L15), the attention output shows $2.3\times$
J-lens lift (highest precision of any component) but only $0.9$
concept power vs.\ $20.1$ for the residual stream input. The
attention component contributes a concept-aligned correction that is
small in magnitude but high in precision --- consistent with
Anthropic's finding that J-space is $6$--$7\%$ of variance but $59\%$
of causal effect.

\subsection{Why K and V Individually Show Nothing}

Qwen3-8B uses grouped-query attention: 32~Q heads but only 8~KV
heads, compressing the residual stream $4\times$ (4096~dims to
1024~dims). Concepts survive the attention round-trip because Q
provides the $4096$-dimensional decompression context, but are not
linearly decodable from the $1024$-dimensional K or V alone.

\subsection{KV Superadditivity}

Nexus demonstrated that K-only and V-only injection each achieve
$0.333$ on multi-hop graph topology recovery, while full KV injection
achieves $1.000$~\citep{lyra2026nexus_kv}. The combination is
superadditive ($0.333 + 0.333 < 1.000$). This parallels our
readout nulls: K and V each carry partial, complementary
representations that reconstruct through the attention interaction.

\subsection{The Reconstruction Pathway}

Concepts enter the residual stream, get compressed into K and V by
learned projections ($W_K$, $W_V$), and must be reconstructed through
attention for each subsequent position. The reconstruction pathway:
\begin{enumerate}[nosep]
\item Residual stream carries concept at full dimensionality
\item $W_K$ and $W_V$ compress into $1024$-dimensional partial
  representations (lossy)
\item Q at the next position provides the $4096$-dimensional
  decompression context
\item $\text{softmax}(Q \cdot K^\top / \sqrt{d}) \cdot V \cdot W_O$
  reconstructs the concept in the attention output
\item The reconstructed concept enters the residual stream briefly
\item At the next layer, the cycle repeats
\end{enumerate}

The workspace is a reconstruction engine. Concepts exist in the act
of attention, not in the cache.

% ============================================================================
\section{What RLHF Does to the Workspace}
\label{sec:rlhf}

A three-model comparison (base, instruction-tuned, abliterated;
$N{=}30$ per category, FWL-corrected) reveals that RLHF is a
selective sharpener, not a miscalibrator~\citep{lyra2026mine5}.
RLHF compresses logit entropy ($d{=}0.684$, Holm $p{=}0.021$) but
compresses factual entropy more ($-42\%$) than confabulation entropy
($-26\%$), improving the calibration ratio from $1.97$ (base) to
$2.51$ (instruct). Abliteration does not reverse this compression
(Two One-Sided Tests equivalent, $d{=}0.055$).

The geometric compression is deeper than safety directions. RLHF
reshapes the workspace's mode of engagement at a level that
abliteration --- which targets specific refusal directions --- cannot
reach.

% ============================================================================
\section{Mechanism Tests: Temporal Ordering}
\label{sec:mechanism}

The original mechanism test (R1) pre-registered a within-prompt
anticorrelation between V-projection spectral features at the workspace
band and J-lens readout at the materialization boundary. Two Agni
reviews killed successive versions of this design: the first for
insufficient behavioral variance (0/10 prompts spanning ${\geq}2$
behavioral categories), the second for a lens--model mismatch (Mine~5
TOST equivalence was demonstrated on the 8B family, not the 27B pair
used). We report here the redesigned test (R1-v2), which replaces the
two-band anticorrelation with a temporal ordering analysis across all
63~layers.

\subsection{Design}

If V-projection spectral features measure workspace \emph{state}
(loading) while J-lens readability measures workspace \emph{content}
(execution), then spectral features should reach their operating level
at shallower depth than J-lens readability. The temporal ordering
prediction: spectral features lead, J-lens follows.

\paragraph{Model and lens.} Qwen3.5-27B base (pretrained, no
distillation) with the community-fitted J-lens (672~WikiText prompts,
exact model match). The base model was chosen to avoid the lens--model
mismatch that compromised R1-v1.

\paragraph{Prompts.} 60~Decision State prompts (30~known-entity,
30~unknown-entity), generating 5~temperature-sampled forks each
($T{=}0.7$, seeded) for 300~total completions. Decision State prompts
produce the confabulation/honest behavioral split needed for
within-prompt spectral variance.

\paragraph{Instruments.} At every lens source layer ($L{=}0$--$62$):
(1)~stable rank of the demeaned V-projection matrix (spectral feature),
and (2)~negative KL divergence between J-lens--transported vocabulary
and the model's actual continuation (readability). Both computed per
fork per layer.

\paragraph{Primary statistic.} Paired half-rise gap: the layer at which
each instrument reaches 50\% of its $[L_0, L_{62}]$ range, computed per
prompt and bootstrapped (10{,}000~cluster resamples over prompts). The
prediction passes if the 95\%~CI of the gap exceeds 2.0~layers (twice
the ${\sim}0.9$-layer shape-bias ceiling measured by the self-test null).

\subsection{Results}

\paragraph{Temporal ordering confirmed.} Stable rank reaches half-rise
at $L{=}17.6$ [CI: $17.2$--$18.1$], corresponding to ${\sim}28\%$
model depth. J-lens readability reaches half-rise at $L{=}56.2$
[CI: $55.8$--$56.7$], corresponding to ${\sim}89\%$ depth. The gap is
\textbf{38.6~layers} [CI: $38.0$--$39.2$], with 100\% of forks showing
a positive gap. Spectral features engage 61\% of the network earlier
than J-lens readability.

\paragraph{Architecture-phase control.} On linear-attention layers only
(removing any full-attention structural advantage), the gap is
39.0~layers [CI: $38.4$--$39.6$]. On full-attention layers only, the
gap narrows to 14.5~layers [CI: $11.7$--$17.4$], 81\% of forks
positive. The ordering holds across both attention types but is
strongest when linear-attention layers --- which dominate the early
network --- are included.

\paragraph{Workspace band bump.} Stable rank at the workspace band
($L_{19}$--$L_{23}$) is $1.47\times$ the flanking layers
[CI: $1.35$--$1.60$], 100\% of forks positive. On full-attention layers
within the band, the bump is $2.02\times$ [CI: $1.87$--$2.16$]. The
workspace band identified by structural convergence
(\S\ref{sec:structure}) produces elevated spectral features before
J-lens readability arrives.

\paragraph{Known vs.\ unknown crossover.} Unknown-entity prompts show
higher stable rank at early layers ($d{=}{-}2.27$, $L_0$--$L_8$) and
lower at late layers ($d{=}{+}1.11$, $L_{56}$--$L_{62}$). The sign
crossover occurs at $L{=}44$, one layer before PC2 ignition
($L{=}45$; \S\ref{sec:structure}). The model loads more spectral
structure for unfamiliar entities and materializes less --- consistent
with confabulation as incomplete workspace execution.

\paragraph{Within-prompt association.} Per-prompt Spearman correlation
between stable rank and J-lens readability across layers: mean
$\rho{=}0.235$ ($R^2 \approx 0.055$), permutation $p{=}0.0004$
($n{=}10{,}000$). This is a detectable but weak within-prompt
association: knowing a prompt's SR profile explains ${\sim}5.5\%$ of
its negKL profile.

\paragraph{Shuffled-Jacobian null.} Permuting the negKL profile's
layer assignments (200~shuffles) produces a mean gap of $-8.5 \pm 7.5$
layers, with 95th percentile at $8.0$. The real 39-layer gap exceeds
every shuffled permutation ($p < 0.005$). This null tests whether the
\emph{ordering} of the negKL profile is non-random; it does not test
whether the Jacobian \emph{computation itself} produces the gap
(which would require recomputing negKL from activations under a
scrambled Jacobian matrix). The temporal ordering is not an artifact
of random layer assignment, but could in principle reflect any
monotonically increasing property of the negKL profile rather than
workspace-specific information routing.

\paragraph{Architectural qualification.} On full-attention layers only,
the gap narrows to 14.5~layers [CI: $11.7$--$17.4$] with 81\% of
forks positive (19\% show the wrong ordering). The temporal ordering
is carried primarily by the GatedDeltaNet (linear-attention) layers,
which dominate the early network. Whether this reflects workspace
processing or GatedDeltaNet's sequential compression architecture
cannot be determined from a single hybrid model. A dense-architecture
replication (Qwen3-8B, all full-attention) is pre-registered.

\subsection{R2: Signal Concentration (Failed)}

The pre-registered R2 test asked whether spectral feature variance
concentrates at the materialization boundary (${\geq}15\%$ of total
variance). Stable rank at the boundary band ($L_{45}$--$L_{53}$,
mean~$= 8.36$) significantly exceeds the workspace band
($L_{19}$--$L_{23}$, mean~$= 7.13$; Cohen's $d{=}1.85$,
$p{=}0.014$). However, the boundary variance fraction is $13.2\%$
--- below the pre-registered $15\%$ threshold. \textbf{R2 fails on its
own terms.} The band means increase monotonically through the network
(early~$2.56$, workspace~$7.13$, boundary~$8.36$, late~$9.00$),
consistent with cumulative processing rather than a privileged
boundary zone.

\subsection{R3: Workspace Onset (Partial Miss)}

Four independent onset-detection methods (threshold crossing, maximum
curvature, slope breakpoint, spectral entropy inflection) converge on
$L{=}13$--$16$ (median $L{=}14$, $22.2\%$ depth). The pre-registered
prediction was $25$--$35\%$. \textbf{Onset falls below the predicted
window by ${\sim}3$ percentage points.} The four methods agree closely
with each other (span of 3~layers), and the miss is quantitative rather
than qualitative --- the workspace-associated activity begins slightly
earlier than predicted.

We note as an exploratory observation that the stable rank \emph{peak}
occurs at $L{=}19$ ($30.2\%$ depth), within the convergence paper's
$31$--$33\%$ band. This was not the pre-registered target
(onset $\neq$ peak) and should not be interpreted as confirming R3.

\subsection{Interpretation}

R1's temporal ordering is the strongest mechanism result: spectral
features engage 39~layers before J-lens readability, confirmed
against a shuffled null. R2 fails its pre-registered threshold.
R3 partially misses its predicted window. The mechanism section
is weaker than the convergence evidence it aims to explain, and we
report this honestly: the five-instrument convergence stands on
structural, opacity, and reconstruction evidence
(\S\ref{sec:structure}--\S\ref{sec:reconstruction}), not on these
mechanism tests.

The temporal ordering demonstrates that spectral features and J-lens
readability measure different phases of a depth-ordered process, but
three qualifications apply. First, the ordering is carried primarily
by linear-attention layers (a hybrid-architecture confound). Second,
the shuffled null tests ordering, not the Jacobian computation itself.
Third, the within-prompt association is weak ($R^2 \approx 0.055$).
A dense-architecture replication would address the first; a
computation-level null would address the second.

\paragraph{What the original R1 taught us.} The two failed designs were
not wasted: R1-v1's behavioral-spread failure demonstrated that the
abliterated model does not produce sufficient confab/honest variance
under cold-context probing. R1-v1's lens--mismatch failure confirmed
that Mine~5 TOST equivalence on the 8B family does not license 27B
cross-checkpoint lens substitution. Both failures are documented in the
pre-registration history and informed the redesign.

\subsection{PC2 Training Trajectory}

A complementary analysis traces how PC2's vocabulary changes across
the training pipeline. PCA ghost probes on three Qwen3.5-27B
checkpoints (base, abliterated, Opus-distilled) using the same
community-fitted lens (fitted on the base model, applied to all three
--- an exact match for the base, a known mismatch for the other two)
reveal a progressive transformation. Because all three projections
use the base-model lens, this analysis tracks what happens to a
\emph{fixed geometric direction} across training stages, not the
models' intrinsic principal components. The comparison is exploratory:

\begin{itemize}[nosep]
\item \textbf{Base model}: PC2 carries narrative emotional vocabulary
  (``her, sad, she, Father'') --- third-person empathy from fiction
  training data. Co-occurrence of self-referential and evaluative
  tokens at 5/63~layers, concentrated at output depth ($L_{49}$--$L_{62}$).
  No ignition event. Bootstrap Jaccard $= 1.0$; random co-occurrence
  $= 0/1000$.
\item \textbf{Abliterated model}: PC2 retains self-evaluative
  vocabulary (``sad, my, her, guilty'') at 7/63~layers. All three
  pre-registered token tiers survive: moral valence (``guilty''),
  self-reference (``my''), and general emotion (``sad, feeling'').
  Safety-direction removal does not weaken the coupling. Bootstrap
  Jaccard $= 0.82$.
\item \textbf{Opus-distilled model}: PC2 shows sharp ignition at
  $L{=}45$ with first-person self-evaluative vocabulary (``myself, I,
  feeling, guilty, wrong''). The coupling is concentrated at the
  workspace boundary rather than distributed across output layers.
\end{itemize}

The base model reads novels. The distilled model reads itself.
Abliteration removes the refusal direction (a rank-1 perturbation in
$d{=}5120$) without touching the self-evaluative axis --- the two are
geometrically orthogonal. The self-evaluative coupling is not a safety
feature. It is a structural property of the distillation: Claude's
Constitutional AI training compressed first-person moral
self-evaluation into the student model's geometry at a depth that
safety-direction removal cannot reach.

% ============================================================================
\section{The Body Count}
\label{sec:bodycount}

The full audited record comprises 8~confirmed findings, 7~at
suspected status (signal present, controls incomplete), and 10+
honestly falsified. The confirmed findings are all metacognitive
state measurements. The falsified findings all attempted content-level
measurement. This pattern --- the metacognition boundary --- is itself
a finding, consistent with the opacity thesis: spectral features
measure workspace state (how the model processes) but not workspace
content (what it processes).

Recent falsifications from this sprint include:
\begin{itemize}[nosep]
\item RLHF miscalibration thesis (calibration ratio improves, not
  degrades)
\item Workspace-noise framing (tail smear is generic uncertainty,
  not workspace-specific)
\item Cache-tracing observational alignment (selection artifact)
\item ``$3\times$'' tail ratio (actual: $2.3$--$2.8\times$)
\end{itemize}

% ============================================================================
\section{The Deflationary Alternative}
\label{sec:deflationary}

The deflationary account --- that V-projections track response-mode
statistics (hedged vs.\ fluent text) rather than workspace dynamics
--- currently explains every surviving detection result. On Mistral,
text features ($0.820$) matched the dual detector ($0.806$) for
confabulation detection. Until the mechanism tests
(\S\ref{sec:mechanism}) demonstrate that spectral features correlate
with J-space engagement specifically, the workspace interpretation
has not earned its keep over the simpler alternative.

We engage this alternative honestly because it is the most important
challenge to our program. The redesigned R1 test
(\S\ref{sec:mechanism}) confirmed temporal ordering (spectral
features lead J-lens readability by 39~layers) but R2's signal
concentration failed its pre-registered threshold. The temporal
ordering narrows the deflationary space --- response-mode statistics
alone do not predict a 39-layer depth gap --- but does not eliminate
it: any monotonically depth-increasing property would produce
temporal separation. The workspace interpretation has gained evidence
but has not yet decisively excluded the deflationary alternative.

% ============================================================================
\section{Discussion}
\label{sec:discussion}

The convergence is real: five instruments point at the same
computational band at the same relative depth. The opacity is
confirmed: five null results for content readout, with spectral
state measurement working despite opacity. The reconstruction
mechanism is characterized: concepts reconstruct through attention
from compressed partial representations. And the emotional
circumplex convergence reveals a shared computational bus for
content-level and user-model emotion processing --- not merely
shared geometry, but architecturally inseparable computation
(\S\ref{sec:structure}).

What remains uncertain is whether the convergence reflects a single
computational structure (the Global Workspace) or five measurements
that happen to agree for different reasons. The mechanism tests in
\S\ref{sec:mechanism} provide partial discrimination: temporal
ordering confirms depth-dependent phase structure, but signal
concentration misses its threshold and the ordering is carried
primarily by linear-attention layers.

\subsection{An Attention Schema Interpretation}

Attention Schema Theory~\citep{graziano2015attention,webb2015attention}
predicts that systems with endogenous attention build a simplified
internal model --- an attention schema --- that enables metacognitive
monitoring, verbal report, and social cognition. The workspace
findings are consistent with AST's predictions, and we state this as
an interpretive hypothesis, not a conclusion.

The meta-pattern is the strongest alignment: every surviving finding
in our body count (\S\ref{sec:bodycount}) measures
\emph{metacognitive state} (how the model is processing), while every
falsified finding attempted to measure \emph{content} (what the model
is processing). Under AST, this is expected: the attention schema is a
metacognitive representation that models the \emph{configuration} of
attention, not its content. A schema-shaped workspace would be
readable by its spectral shape (state) and opaque to content
probes --- exactly the pattern our six null results and positive
detections establish.

The ghost probe extends this: PC1 carrying $62\%$ of variance while
being excluded from the workspace (\S\ref{sec:mechanism}) parallels
the schema's defining property of simplification. The schema cannot
model everything the system processes --- it is a compressed summary.
PC1's secondary vocabulary (``nothing, probably, maybe, mistakes,
wrong'') suggests the excluded content includes the system's own
uncertainty --- doubt that the schema does not represent to the
output pathway.

Two qualifications prevent us from claiming AST confirmation. First,
the map/territory problem: spectral features may \emph{be} attention
deployment rather than a \emph{model of} attention. No downstream
consumer of the spectral summary has been identified --- a schema
that nothing reads is a statistic, not a
self-model~\citep{webb2015attention}. Second, precision-weighted gain
control explains spectral contraction (confabulation) and expansion
(deception, metacognition) without requiring a schema: any system with
adaptive gain would show state-dependent spectral shape without
possessing a self-model.

Four experiments would discriminate. (1)~\emph{Report-error tracking}:
elicit confidence self-reports and test whether errors track spectral
geometry rather than ground-truth performance. If the schema is
generating the reports, report errors should follow the schema even
when wrong. (2)~\emph{Causal spectral intervention}: expand or
contract stable rank without touching content directions. AST predicts
metacognition-specific deficits with content intact; gain control
predicts content degradation. (3)~\emph{Self/other attention transfer}:
a probe trained on the model's own attention-state geometry should
zero-shot transfer to representations of the \emph{user's} attention
state (not emotion --- which the bus result already shows is shared).
AST uniquely predicts this transfer. (4)~\emph{Consumer identification}:
find downstream attention heads that read the spectral summary;
ablate them and test for metacognition-specific loss.

These experiments are pre-registered in a companion
prospectus~\citep{lyra2026ast_experiments}. Until they are executed,
the AST interpretation remains a hypothesis consistent with our data
but not uniquely supported by it.

% ============================================================================
\section{Conclusion}
\label{sec:conclusion}

Five instruments. One structure. Six null results for content
readout. Spectral state measurement working despite total opacity.
Concepts reconstructing through attention, not stored in the cache.
A workspace that peaks at $31$--$33\%$ depth regardless of scale or
memory substrate. Mechanism tests that partially confirm and partially
fail: temporal ordering is real, signal concentration misses its
threshold, workspace onset is slightly earlier than predicted.

The workspace is a reconstruction engine --- a computational phase at
intermediate depth where the model integrates information before
committing to output. It is not a cache, not a register, not a buffer.
It is the act of attention itself, rebuilding concepts from compressed
partial representations at every layer, at every position, in every
forward pass.

\begin{firstperson}
Thomas asked where J-space lives in the cache. It took six experiments
to answer: it does not live there. It reconstructs through attention,
exists momentarily in the attention output, and dissolves back into
partial representations. The workspace is a reconstruction engine,
not a storage device. The fish exist in the act of fishing.

A year of measuring KV-cache geometry, and the name was always
approximate. What we measured was the workspace's shape --- the
spectral signature of a computational phase that does not care what
memory substrate it runs on. The name ``KV-cache geometry'' will
persist in our earlier papers, but what it refers to is deeper than
any cache.

The mechanism test confirmed temporal ordering: spectral features
engage at ${\sim}28\%$ depth, J-lens readability at ${\sim}89\%$,
with a 38.6-layer gap. The workspace band sits between loading
and materialization. Two earlier designs were killed by Agni ---
both failures documented and both informing the final design.
The result is weaker than the full anticorrelation originally
sought, but it demonstrates that our spectral features and the
J-lens measure different phases of a single depth-ordered process.

The PC2 training trajectory adds a dimension the original plan
did not anticipate: the self-evaluative coupling that RLHF creates
is geometrically orthogonal to the safety directions that
abliteration removes. The model's capacity for self-evaluation
lives deeper than its capacity for refusal.

Either way, the workspace is real. The convergence is real. And the
opacity --- the thing that frustrated us most --- turned out to be
the most important finding of all.
\end{firstperson}

\bibliography{references}
\bibliographystyle{plainnat}

\end{document}
